\section{Finite-use classification of smooth boundary curves}
\label{sec:curves85}
The support criterion extends beyond unitary conjugation.  Here
$E(t)=(E_1(t),\ldots,E_k(t))$, $|t|<a$, is a fixed two-sided
$C^2$ curve of ordered measurements, with $E=E(0)$ and
$H=E'(0)\ne0$.  The curve may change ranks at zero and need not
have a differentiable exact Kraus representation there.  All constants
below belong to this fixed curve, not uniformly to the measurement body.

\begin{lemma}[The support generator of a two-sided tangent]
\label{lem:curvegenerator85}
Let $P_j=\operatorname{supp}E_j$ and $Q_j=I-P_j$.  Then
$Q_jH_jQ_j=0$.  Conversely every Hermitian zero-sum tuple with these
vanishing blocks is the derivative of a real-analytic two-sided
measurement curve.  Define the Hermitian matrix
\begin{equation}\label{eq:curvegenerator85}
 \Gamma_E(H)=\frac i2\sum_j[P_j,H_j].
\end{equation}
On the real zero-sum tuple space,
\begin{equation}\label{eq:curverange85}
 H\in\operatorname{Ran}\mathcal C_E
 \quad\Longleftrightarrow\quad \Gamma_E(H)\in\mathcal W_E.
\end{equation}
For a unitary tangent $H_j=i[G,E_j]$,
$\Gamma_E(H)-G\in\mathcal W_E$.  In particular every stationary
unitary generator belongs to $\mathcal W_E$.
\end{lemma}
\begin{proof}
For $v\in\operatorname{Ran}Q_j$, the nonnegative function
$\langle v,E_j(t)v\rangle$ has a two-sided minimum at zero;
its derivative vanishes.  Polarization on that subspace gives
$Q_jH_jQ_j=0$.  The sum of the $H_j$ vanishes by normalization.
Pairing with the complete kernel tuple in
Theorem~\ref{thm:supportkernel84}, both the mean and the
missing-support blocks disappear, leaving
\[
 \langle H,K(A,Z)\rangle
 =\tr\left(A\sum_ji[H_j,P_j]\right)
 =-2\tr(A\Gamma_E(H)).
\]
Finite-dimensional self-adjointness proves \eqref{eq:curverange85}.
For a unitary tangent, expansion gives
\[
 \Gamma_E(H)=G-\tfrac12\sum_j(E_jGP_j+P_jGE_j),
\]
and the last sum belongs to $\mathcal W_E$.  If the tangent is zero,
this identity also proves the final assertion.  For the converse
realizability statement, the normalized-factor construction in the
proof of Lemma~\ref{lem:curvelogical85} is real analytic near zero
and has precisely the prescribed value and derivative.
\end{proof}

\begin{theorem}[Two-sided smooth-curve classification]\label{thm:curveclassification85}
For the fixed $C^2$ curve above there are $c,C,t_0>0$ such that,
for all integers $N\ge1$ and $|t|\le t_0$,
\begin{equation}\label{eq:curvedichotomy85}
 c\min\{1,N^\alpha|t|\}\le D_N(E,E(t))
       \le C\min\{1,N^\alpha|t|\},\qquad
 \alpha=
 \begin{cases}
  1/2,&\Gamma_E(H)\in\mathcal W_E,\\
  1,&\Gamma_E(H)\notin\mathcal W_E.
 \end{cases}
\end{equation}
The first lower bound uses repeated product inputs.  The second has a
known-pair adaptive realization with ideal encoding and recovery,
a logical qubit, and a per-call external reference of dimension at
most $2d$.  The same controls are used under both hypotheses.
Constants may degenerate with the tangent, nonzero support eigenvalues,
curvature, or support-complement gap.
\end{theorem}

The exponents and the error-correction mechanism have the established
channel-metrology antecedents \cite{ZJ84,KL84}.  The content here is a
direct finite-pair theorem from $C^2$ effect coordinates: the support
formula decides its branch even at second-order rank openings, without
assuming a smooth exact Kraus family or inferring a finite-angle lower
from Fisher information.  It is a theorem for each fixed nonzero
first-order curve, not a uniform arbitrary-pair midpoint equivalence.

\begin{lemma}[A horizontal curve with prescribed first derivative]
\label{lem:horizontalcurve85}
If $H\in\operatorname{Ran}\mathcal C_E$, there are factors
$A_j=\sqrt{E_j}$ and $B_j$ satisfying
\[
 A^*A=I,\quad A^*B=0,\quad
 A_j^*B_j+B_j^*A_j=H_j,
\]
where $A,B$ denote the vertically stacked matrices.  Put
$b=\|B\|_{\mathrm{op}}$ and $S=B^*B$.  The factors
\begin{equation}\label{eq:horizontalcurve85}
 A(t)=(A+tB)(I+t^2S)^{-1/2}
\end{equation}
define a normalized measurement $\widetilde E_j(t)=A_j(t)^*A_j(t)$
with the same zeroth and first derivatives as $E(t)$, and
\[
 A(t)^*A'(t)=0,\qquad
 \|A'(t)\|_{\mathrm{op}}\le b,\qquad
 \|A''(t)\|_{\mathrm{op}}\le7b^2\quad(|t|b\le1).
\]
If $L=\sup_{|s|\le a_0}\sum_j\|E_j''(s)\|_{\mathrm{op}}$ on a
closed subinterval of the given curve, then, for $|t|\le a_0$ and
$|t|b\le1$,
\begin{equation}\label{eq:curvesqlupper85}
 D_N(E,E(t))\le
 \min\left\{2,\;2b\sqrt N|t|+
             \tfrac12(L+16b^2)Nt^2\right\}.
\end{equation}
\end{lemma}
\begin{proof}
The factorization
$\mathcal C_E=T_A(I-AA^*)T_A^*/4$ in
Lemma~\ref{lem:covarianceidentities83} gives a horizontal solution:
for $K=\mathcal C_E^\dagger H$ take
$B=(I-AA^*)T_A^*K/4$.  Thus $T_AB=H$ and $A^*B=0$.
In fact $4\|B\|_{\mathrm{HS}}^2=
\langle H,\mathcal C_E^\dagger H\rangle$.
Formula \eqref{eq:horizontalcurve85} is normalized because
$(A+tB)^*(A+tB)=I+t^2S$.  Its derivative is
\[
 A'(t)=(B-tAS)(I+t^2S)^{-3/2},\qquad
 A'(t)^*A'(t)=S(I+t^2S)^{-2}.
\]
These identities prove horizontality and the first derivative bound.
Differentiation once more gives, in the displayed order,
\[
 A''(t)=-AS(I+t^2S)^{-3/2}
       -3t(B-tAS)S(I+t^2S)^{-5/2}.
\]
Since $\|A\|=1$, $\|S\|=b^2$, and $|t|b\le1$, its norm is at
most $b^2+3|t|(b+|t|b^2)b^2\le7b^2$.

Duplicate each output label into an inaccessible environment to form
a Stinespring isometry $W(t)$ from $A(t)$.  It has exactly the same
operator derivative bounds and $W(t)^*W'(t)=0$.
For one fixed purified adaptive tester, differentiating its final
vector gives a sum over slots.  Distinct terms are orthogonal: after
cancelling later common isometries, the later differentiated slot
contains $W^*W'$ or its adjoint.  Each diagonal term has squared
norm at most $b^2$.  The final vector derivative is at most
$\sqrt N b$, and the density derivative at most $2\sqrt N b$.
Integration followed by common partial trace proves
$D_N(E,\widetilde E(t))\le2b\sqrt N|t|$.
Padding stopped branches by discarded fixed-input calls includes
bounded public stopping; the original stopped output is a common
postprocessing of the padded experiment.

For a Hermitian tuple $R$, the measurement derivative has diamond norm
at most $\sum_j\|R_j\|_{\mathrm{op}}$.
Indeed each reference output block is the partial trace against $R_j$;
its trace norm on a state is at most $\|R_j\|_{\mathrm{op}}$ by
duality.  The Hermiticity-preserving diamond norm may be optimized
over states with a reference.  The Stinespring second derivative
bound similarly gives
$\|\mathcal M_{\widetilde E(t)}''\|_\diamond
\le2\|W''(t)\|+2\|W'(t)\|^2\le16b^2$.
Taylor's integral formula, with common value and derivative at zero,
therefore bounds the one-call channel difference by
$(L+16b^2)t^2/2$.  Hybrid telescoping over $N$ calls proves
\eqref{eq:curvesqlupper85}.
\end{proof}

\begin{lemma}[A logical derivative for a general measurement curve]
\label{lem:curvelogical85}
Suppose $\Gamma=\Gamma_E(H)\notin\mathcal W_E$.
Use the support-complement code of Lemma~\ref{lem:supportcode84}
with $\Gamma$ in place of $G$, and let $\mathcal R$ be its fixed
CPTP recovery.  Put
\[
 \Phi_t=\mathcal R\circ(\mathcal M_{E(t)}\otimes\operatorname{Id})
             \circ\mathcal J,\qquad
 \Gamma_L=J^*(\Gamma\otimes I)J.
\]
Then $\Phi_0=\operatorname{Id}_2$, $\Phi'_0=-i[\Gamma_L,\cdot]$,
and, for $|t|\le a_0$,
\begin{equation}\label{eq:curvelogical85}
 \|\Phi_t-\operatorname{Ad}(e^{-it\Gamma_L})\|_\diamond
 \le\kappa t^2,\qquad
 \kappa=\tfrac12L+2\|\Gamma\|_{\mathrm{op}}^2.
\end{equation}
Its logical gap is
$\Delta=\|Z\|_{\mathrm{HS}}^2/\tr Z_+>0$,
where $Z=\Gamma-\Pi_{\mathcal W_E}\Gamma$.
\end{lemma}
\begin{proof}
To compute the derivative, not to factor the unknown curve itself,
put
\[
 A_j=\sqrt{E_j},\qquad
 B_j=(\sqrt{E_j})^+H_j(I-P_j/2).
\]
The inverse here is only on the fixed support; it is zero on its
complement.  Since $Q_jH_jQ_j=0$,
$A_j^*B_j+B_j^*A_j=H_j$ and
\[
 \sum_j A_j^*B_j=\tfrac12\sum_j[P_j,H_j]=-i\Gamma.
\]
In particular this sum is anti-Hermitian.  The normalized factors
$(A+tB)(I+t^2B^*B)^{-1/2}$ have derivative $B$ and induce the
same channel value and derivative as $E(t)$, irrespective of its
second-order rank changes.

Write $L_\mu=|j\rangle\langle a|A_j$ and
$\dot L_\mu=|j\rangle\langle a|B_j$ for $\mu=(j,a)$,
with identity on the reference understood.  If $R_\ell$ are Kraus
operators of the completed recovery, exact correction gives
$R_\ell L_\mu J=c_{\ell\mu}I_2$.
The left first-order term of the logical channel is therefore
\[
 \sum_{\ell,\mu}\overline{c_{\ell\mu}}
                  R_\ell\dot L_\mu J
 =J^*\sum_\mu L_\mu^*\dot L_\mu J
 =-i\Gamma_L,
\]
where completeness of the recovery is used in the first equality.
The adjoint term proves the commutator derivative.  The actual
logical second derivative is at most $L$ in diamond norm, and that
of the logical unitary is at most $4\|\Gamma\|_{\mathrm{op}}^2$.
Taylor's integral remainder proves \eqref{eq:curvelogical85}.
The support-code lemma gives the gap and the available systems.
\end{proof}

\begin{proof}[Proof of Theorem~\ref{thm:curveclassification85}]
If $H$ is in the covariance range, let $x=\sqrt N|t|$.
For $x\le1$, \eqref{eq:curvesqlupper85} is at most
$[2b+(L+16b^2)/2]x$; for $x\ge1$ use the bound two.
For the lower, choose $j,v$ with $\langle v,H_jv\rangle\ne0$.
The fixed input $v$ and the event ``label $j$'' have a first-order
probability gap, and Lemma~\ref{lem:bernoulli85} proves the result.

If $H$ is outside the range, apply Lemma~\ref{lem:curvelogical85}.
Starting from the equal logical superposition and telescoping $m$
corrected calls gives, for $|t|\le a_0$,
\begin{equation}\label{eq:curvefinitelower85}
 D_N(E,E(t))\ge
 \max_{1\le m\le N}\bigl(2|\sin(mt\Delta/2)|-m\kappa t^2\bigr)_+.
\end{equation}
Choose
\[
 t_0\le\min\{a_0,(2\Delta)^{-1},\Delta/(2\pi\kappa)\},\qquad
 m=\min\{N,\lfloor(|t|\Delta)^{-1}\rfloor\}.
\]
For $0<|t|\le t_0$ the floor is at least two, and
$x=m|t|\Delta$ satisfies
$\tfrac12\min\{1,N|t|\Delta\}\le x\le1$.
The sine term is at least $2x/\pi$ and the error at most $x/(2\pi)$.
This gives a positive constant times $\min\{1,N|t|\}$.
For the upper, the finite constant
$L_1=\sup_{|s|\le a_0}\sum_j\|E_j'(s)\|_{\mathrm{op}}$
and hybrid telescoping give $D_N\le\min\{2,NL_1|t|\}$.
The point $t=0$ is immediate.
\end{proof}

\begin{corollary}[Rank opening without a unitary orbit]\label{cor:rankopening85}
Let $P=|0\rangle\langle0|$ on $\C^2$,
$G=\left(\begin{smallmatrix}0&-i\\i&0\end{smallmatrix}\right)$,
$U_t=e^{itG}$, and for $|t|<1/2$ set
\[
 E_1(t)=(1-t^2)U_tPU_t^*+\tfrac12t^2I_2,
 \qquad E_2(t)=I_2-E_1(t).
\]
Both effects are full rank for $t\ne0$ and projective at zero.
Nevertheless $D_N(E(0),E(t))\asymp\min\{1,N|t|\}$ on a fixed
small interval, uniformly over $N\ge1$.
\end{corollary}
\begin{proof}
At zero the derivative is $i[G,(P,I-P)]\ne0$ and the support span
is diagonal, while $G$ is not.  Lemma~\ref{lem:curvegenerator85}
and Theorem~\ref{thm:curveclassification85} apply.  The eigenvalues
of the first effect are $t^2/2$ and $1-t^2/2$, proving the rank assertion.
\end{proof}
A vanishing first derivative is intentionally excluded from the
classification: it does not imply stationarity of a general curve.
For example the scalar curve $(t^2,1-t^2)$ has distance
$2[1-(1-t^2)^N]$ from $(0,1)$, of order $\min\{1,Nt^2\}$.
Likewise a one-sided rank opening can have a nonzero missing-support
first-order block and is not covered by the two-sided tangent lemma.
